\section{Durable Trust: el efecto compuesto de Reliability, Security y la recuperación transparente}

Los primeros clientes de una identity startup tienen que confiar en que un equipo pequeño puede proteger la puerta de entrada de la organización, sin paros frecuentes por mantenimiento, y en que asumirá su responsabilidad cuando ocurra un incidente. La experiencia de McKinnon en Salesforce aporta la credibility inicial, pero la confianza a largo plazo solo se acumula con production evidence: servicio estable, change correctos, incident response honesta, customer guidance e inversión continua en la plataforma.

\subsection{Trust Flywheel}

Un servicio confiable lleva a los clientes a ampliar su deployment. Más aplicaciones y más usuarios aportan una integration/context más rica, pero también amplían la responsabilidad. La security evidence y un RCA transparente ayudan al cliente a evaluar el residual risk. A su vez, la confianza del cliente trae renovaciones, reference y una adopción más profunda de la plataforma. El volante también puede girar en sentido contrario: un poorly handled incident no solo afecta el SLA del periodo, sino que lleva al cliente a reducir el alcance de su integración, lo que debilita el context futuro y el valor del producto.

\lecturefigure{15-trust-flywheel.png}{La adopción de una plataforma de identidad genera un efecto compuesto mediante la confiabilidad, la evidencia de seguridad, la confianza del cliente y la profundidad de la plataforma.}{Subtítulos locales 27:30--35:20; redibujo conceptual.}

\begin{knowledgebox}{Lectura de la figura: Trust es evidencia acumulada, no una promesa de cero incidentes}
El reliable service demuestra que se puede depender de la front door; la security evidence muestra el control y la capacidad de aprendizaje; la customer trust se refleja en implementaciones más amplias, reference y renewal; y la platform depth, a su vez, mejora la integration y el policy context. Ningún sistema grande puede prometer con honestidad que nunca tendrá incidentes, pero sí puede comprometerse a detectarlos rápido, limitar el blast radius, publicar hechos sobre los que se pueda actuar y completar una remediation verificable.
\end{knowledgebox}

\teachervoice{El ponente explicó que los primeros clientes no solo compran una funcionalidad, también evalúan «si esta persona se irá en cuanto cierre la venta y si es capaz de mantener estable el servicio». Su experiencia previa en ingeniería le dio la primera ronda de confianza, pero cada outage, breach y recuperación posterior vuelve a fijar el valor de esa confianza.}

\subsection{Resumen de la sección}

La identity trust se construye con evidencia operativa de largo plazo. Reliability, security, transparency y customer outcome se refuerzan entre sí; el eslogan de marca solo puede amplificar la evidencia que ya existe, no sustituirla.
